\section{Actual deletion of an interior identity empty step at every depth}

Keep the original marked record and every original indexed member. Delete
only the actual identity empty blowup between the original head and suffix;
the entire original suffix and all its centres remain literal. The
ordered-boundary route of \cite{wlodarczyk05-hironaka,bierstone-milman08-functoriality}
is adapted by deriving original prefix canonicality during deletion, before
using the established whole-value and dominance transport.

\begin{theorem}[Whole original canonical histories reflect interior identity empty deletion]
\label{hir:original-recursive-interior-empty-canonical-deletion-iff}
\source{wlodarczyk05-hironaka,bierstone-milman08-functoriality}
\uses{hir:original-recursive-entire-family-interior-history,hir:original-recursive-entire-step-interior-history,hir:original-recursive-interior-spectrum-terminal-history}
\lean{Hironaka.HironakaOriginalRecursiveEmptyDeletion.original_recursive_family_empty_deletion_isCanonical_iff,Hironaka.HironakaOriginalRecursiveEmptyDeletion.original_recursive_step_empty_deletion_isCanonical_iff}
\leanok
For EVERY original family depth, original bound-$k$ record $M$, original
head $H$ and entire suffix $S$, the literal inserted history $H\epsilon S$
is canonical for the actual $R_k$ if and only if $HS$ is canonical.
The same statement holds for the actual Step One $R_k$ at original
bound-$(k+1)$ maximal-order inputs. In the deletion direction ONLY the actual
inserted whole canonical history is assumed; no deleted canonicality or
value comparison certificate is supplied.
\end{theorem}
\begin{proof}
Insertion follows the established original-input whole canonical-history
producer. For deletion induct on the SAME suffix $S$. If it is empty,
canonicality of the last identity empty step contains original head
canonicality directly. At a next actual step, inserted whole canonicality
supplies inserted prefix canonicality and the actual centre condition.
Apply the induction hypothesis FIRST to construct deleted OLD prefix
canonicality. That genuine prefix now satisfies the hypotheses of the
established actual dominance dictionary. If the centre is empty it remains
empty. Otherwise it equals the entire inserted prefix dominance locus;
the proved dictionary identifies that locus with the entire deleted prefix
dominance locus. Thus the literal OLD next step is canonical. This constructs
deleted whole canonicality internally and proves both equivalences. No
maximum existence, branch, level or lower-order condition is required.
\end{proof}

\begin{theorem}[Actual inserted canonical input internally computes the deleted whole values]
\label{hir:original-recursive-interior-empty-deletion-entire-values}
\source{wlodarczyk05-hironaka,bierstone-milman08-functoriality}
\uses{hir:original-recursive-interior-empty-canonical-deletion-iff,hir:original-recursive-entire-family-interior-history,hir:original-recursive-entire-step-interior-history,hir:original-interior-empty-geometric-dictionary}
\lean{Hironaka.HironakaOriginalRecursiveEmptyDeletion.original_recursive_family_empty_deletion_entire_values,Hironaka.HironakaOriginalRecursiveEmptyDeletion.original_recursive_step_empty_deletion_entire_values}
\leanok
ONLY original domain data and actual INSERTED whole canonicality internally
supply the deleted whole canonical history. At every actual inserted
cosupport point, the entire inserted family value equals $H_{k,j}$ of the
literal deleted whole value, and the inserted Step One value equals
$S_{k,j}$ of the deleted value, where $j=|M.E|+|H|$. All true levels and
orders and empty suffixes are included; no OLD canonicality, cosupport or
value certificate is a caller premise.
\end{theorem}
\begin{proof}
The first theorem constructs deleted canonicality from the actual inserted
canonical history. The original invariant then supplies genuine testness
of the deleted whole history and original head. The actual geometric
dictionary gives equality of inserted and deleted transformed cosupport.
An inserted cosupport point is therefore internally a genuine deleted
cosupport point. Apply the established ENTIRE original value formula to
this internally supplied OLD canonical history and point. Its SAME strict
chronological map gives the stated complete value without selecting a new
assignment history or rebasing the original marked record.
\end{proof}

\begin{theorem}[Whole canonical terminal histories reflect the same actual deletion]
\label{hir:original-recursive-interior-empty-terminal-deletion-iff}
\source{wlodarczyk05-hironaka,bierstone-milman08-functoriality}
\uses{hir:original-recursive-interior-empty-canonical-deletion-iff,hir:original-recursive-interior-spectrum-terminal-history,hir:original-interior-empty-geometric-dictionary}
\lean{Hironaka.HironakaOriginalRecursiveEmptyDeletion.original_recursive_family_empty_deletion_canonical_resolution_iff,Hironaka.HironakaOriginalRecursiveEmptyDeletion.original_recursive_step_empty_deletion_canonical_resolution_iff}
\leanok
At EVERY family and Step One depth, the literal inserted history is both
canonical and a terminal geometric resolution if and only if the literal
deleted history has these properties. No terminal or canonicality
certificate for the deleted history is supplied in the deletion direction.
\end{theorem}
\begin{proof}
For insertion use the established whole canonical terminal producer.
For deletion the first theorem internally gives deleted canonicality.
The genuine original invariant supplies actual testness of that deleted
history and head. Original transformed cosupport equality then transfers
inserted terminal emptiness to the deleted history. Together these give
the true deleted geometric resolution and canonicality. This transports
an actual given history; arbitrary empty biholomorphic maps, noncompact
canonical scheduling and global resolution gluing remain distinct source
obligations.
\end{proof}
